\subsection{Activity on the wrong protein: pseudoenzymes, complex subunits and the missing NOTs}\label{sec:dd-pseudo}
\textbf{Claim.} Automated channels assign a family's or a complex's catalytic activity to proteins
that do not perform it: fold-retaining but catalytically dead pseudoenzymes, and structural,
accessory or assembly subunits of an enzyme complex. Only a NOT annotation stops re-propagation.
\textbf{Projects.} \proj{PSEUDOENZYMES.md}, \proj{CONTESTED\_FUNCTION.md}: tardigrade Cu/Zn-SOD paralogs,
Epe1, EryCII, PqsB, ActI-ORF2; \textit{C.\ elegans} LYS-7 with lysozyme activity from its domain.
\proj{PROTEIN\_COMPLEX\_FUNCTIONS.md}, \proj{BGC.md}, \proj{SDH\_GP2TERM\_CONTRIBUTES\_TO.md}: catalytic vs
structural proteasome subunits; obligate heterodimers in natural-product clusters; succinate
dehydrogenase subunits moved to \texttt{contributes\_to}. \proj{COMPLEXPORTAL.md}: subunit rows mostly
sound (\accComplexPortal\% accepted here), but failures include the inhibited subunit annotated as the
inhibitor and promiscuous subunits inheriting every complex's terms.
\proj{TOP\_NOTS.md}: 250 NOT candidates mined from removal reasons across 19 species.
\proj{UNIPROT\_CAUTION\_NOTE.md}, \proj{NOT\_ANNOTATION\_USAGE.md}, \proj{IRD\_EVIDENCE.md}: IRD losses never
reach a protein as a NOT.
\textbf{Method highlight.} Catalytic-residue conservation, M-CSA, HMM scoring, structure.
\textbf{Upstream.} \todo{NOTs filed / accepted.}
